\section{Моделирование методом динамики частиц. Молекулярная динамика}
\label{sec:particle-molecular-dynamics}

\subsection{Дискретная модель вещества}

В методе динамики частиц сплошное или атомистическое тело представляется
системой $N$ взаимодействующих материальных частиц. Состояние системы задаётся
координатами и скоростями
\[
    \mathbf r_i(t),\qquad \mathbf v_i(t)=\dot{\mathbf r}_i(t),
    \qquad i=1,\ldots,N,
\]
а также массами $m_i$. В молекулярной динамике частицами обычно являются атомы
или молекулы, а взаимодействие между ними задаётся межатомным потенциалом.

Такой подход принципиально отличается от континуального: вместо полей
перемещений, напряжений и деформаций непосредственно вычисляются траектории
отдельных частиц. Макроскопические величины затем получают усреднением по
частицам и, при необходимости, по ансамблю реализаций.

\subsection{Уравнения движения}

Основу молекулярной динамики составляют уравнения Ньютона
\[
    m_i\ddot{\mathbf r}_i=\mathbf F_i,
    \qquad i=1,\ldots,N.
\]
Если взаимодействие консервативно и потенциальная энергия системы равна
$U(\mathbf r_1,\ldots,\mathbf r_N)$, то
\[
    \mathbf F_i=-\nabla_{\mathbf r_i}U.
\]
Для парного взаимодействия
\[
    U=\sum_{i<j}\varphi(r_{ij}),
    \qquad
    r_{ij}=|\mathbf r_i-\mathbf r_j|,
\]
и сила, действующая на частицу $i$ со стороны $j$, направлена вдоль линии,
соединяющей частицы:
\[
    \mathbf F_{ij}
    =-\varphi'(r_{ij})
      \frac{\mathbf r_i-\mathbf r_j}{r_{ij}}.
\]

Для кристаллической решётки уравнения часто записывают через перемещения
$\mathbf u_i$ частиц из положений равновесия. В гармоническом приближении они
образуют линейную систему вида
\[
    \ddot{\mathbf u}=L\mathbf u,
\]
где $L$ --- разностный оператор, определяемый массами, геометрией решётки и
жёсткостями связей.

\subsection{Потенциалы межчастичного взаимодействия}

Качество молекулярно-динамической модели в значительной степени определяется
выбранным потенциалом. Простейшее гармоническое взаимодействие вблизи
равновесия имеет квадратичный вид
\[
    \varphi(r)\approx
    \frac{k}{2}(r-a)^2,
\]
где $a$ --- равновесное расстояние.

Для моделирования нелинейного взаимодействия в кристаллах часто используют
потенциал Леннарда--Джонса. В обозначениях пособия СПбПУ он записывается как
\[
    \varphi(r)=\varepsilon
    \left[\left(\frac{a}{r}\right)^{12}
          -2\left(\frac{a}{r}\right)^6\right],
\]
где $\varepsilon$ задаёт характерную энергию связи, а $a$ --- равновесное
расстояние. Член $r^{-12}$ отвечает за быстрое отталкивание на малых расстояниях,
а член $r^{-6}$ --- за притяжение.

В простых моделях кристаллов можно учитывать только ближайших соседей. Для
больших систем с короткодействующими потенциалами на практике также ограничивают
расчёт взаимодействий некоторым радиусом, чтобы не перебирать все пары частиц.

\subsection{Начальные и граничные условия}

Для задачи Коши задаются положения и скорости всех частиц:
\[
    \mathbf r_i(0)=\mathbf r_i^0,
    \qquad
    \mathbf v_i(0)=\mathbf v_i^0.
\]
При моделировании теплового движения начальные скорости могут рассматриваться
как случайные величины с заданными статистическими характеристиками.

Для моделирования объёмного кристалла широко применяются периодические
граничные условия. В одномерной решётке условия Борна--Кармана имеют вид
\[
    \mathbf u_{i+N}=\mathbf u_i,
\]
то есть конечная расчётная ячейка рассматривается как периодически повторяющаяся
часть бесконечного кристалла. Это уменьшает влияние искусственных внешних
границ.

\subsection{Численное интегрирование. Метод Верле}

Уравнения движения интегрируются по времени с малым шагом $\Delta t$. В
молекулярной динамике особенно распространены схемы Верле, так как они просты,
имеют второй порядок по координатам и хорошо приспособлены к длительному
интегрированию механических систем.

В пособии СПбПУ метод Верле указан как используемый численный интегратор; ниже для полноты приведена стандартная запись позиционного алгоритма Верле:
\[
    \mathbf r_i^{n+1}
    =2\mathbf r_i^n-\mathbf r_i^{n-1}
     +\Delta t^2\,\mathbf a_i^n,
    \qquad
    \mathbf a_i^n=\frac{\mathbf F_i^n}{m_i}.
\]
Если нужны скорости на каждом шаге, удобно использовать вариант
\emph{velocity Verlet}:
\[
\begin{aligned}
    \mathbf v_i^{n+1/2}&=\mathbf v_i^n
        +\frac{\Delta t}{2}\mathbf a_i^n,\\
    \mathbf r_i^{n+1}&=\mathbf r_i^n
        +\Delta t\,\mathbf v_i^{n+1/2},\\
    \mathbf v_i^{n+1}&=\mathbf v_i^{n+1/2}
        +\frac{\Delta t}{2}\mathbf a_i^{n+1}.
\end{aligned}
\]
На новом положении сначала пересчитываются межчастичные силы, после чего
завершается шаг по скорости. Шаг $\Delta t$ должен быть существенно меньше
характерного периода самых быстрых колебаний системы.

\subsection{Макроскопические характеристики}

По траекториям частиц вычисляют энергию и температурные характеристики.
Кинетическая и потенциальная энергии равны
\[
    K=\sum_{i=1}^{N}\frac{m_i v_i^2}{2},
    \qquad
    U=U(\mathbf r_1,\ldots,\mathbf r_N),
\]
а в изолированной консервативной системе полная энергия
$E=K+U$ в точной динамике сохраняется.

Температура связана со средней кинетической энергией теплового движения. Для
одинаковых частиц в $d$ пространственных измерениях, после исключения движения
центра масс, используется соотношение вида
\[
    \frac{d}{2}k_B T
    =\left\langle\frac{m v^2}{2}\right\rangle.
\]
В неравновесных задачах полезно рассматривать также ковариации скоростей и
тензорную температуру. Именно такие характеристики позволяют описывать
перераспределение энергии между пространственными направлениями.

\subsection{Пример: кристалл Леннарда--Джонса}

В пособии СПбПУ рассматривается треугольная решётка с взаимодействием
частиц по потенциалу Леннарда--Джонса. Учитываются взаимодействия
ближайших соседей, используются периодические граничные условия, а уравнения
динамики численно интегрируются методом Верле. Молекулярно-динамический расчёт
показывает переход системы к стационарному состоянию и выравнивание
кинетической и потенциальной энергий; при малых температурах результаты близки
к гармонической модели, а с ростом нелинейности переходный процесс изменяется.

\subsection{Преимущества и ограничения метода}

Преимущества молекулярной динамики:
\begin{itemize}
    \item прямое описание микроскопического движения и нелинейных процессов;
    \item возможность исследовать неравновесные процессы и получать
    макроскопические характеристики из микродинамики;
    \item отсутствие необходимости заранее задавать континуальные
    определяющие соотношения материала.
\end{itemize}
Основные ограничения --- высокая вычислительная стоимость при большом числе
частиц, малые допустимые шаги по времени и сильная зависимость результата от
адекватности межчастичного потенциала. Кроме того, для тепловых задач часто
нужно статистическое усреднение по частицам или реализациям.

\subsection{Что важно сказать на экзамене}

Главная схема ответа:
\[
    \text{частицы и потенциал}
    \;\longrightarrow\;
    \mathbf F_i=-\nabla_i U
    \;\longrightarrow\;
    m_i\ddot{\mathbf r}_i=\mathbf F_i
    \;\longrightarrow\;
    \text{интегрирование методом Верле}.
\]
Нужно обязательно упомянуть начальные условия, периодические граничные условия
и привести пример потенциала Леннарда--Джонса. После решения траекторий из
скоростей и координат вычисляют энергию, температуру и другие макроскопические
характеристики.

\newpage
